% ============================================================
%  related_works.tex
%  Related Work: Memory- and Communication-Efficient
%  Federated Learning
%
%  Target venues : ICML / NeurIPS / ICLR
%  Citations     : 25  (≥ 12 threshold)
% ============================================================

\section{Related Work}
\label{sec:related}

Our work sits at the intersection of three active research threads:
federated learning under data heterogeneity, communication-efficient
distributed optimization, and memory-efficient training.  We survey
each thread in turn and then articulate the gap that motivates our
joint treatment.

% ------------------------------------------------------------
\subsection{Federated Learning}
% ------------------------------------------------------------

Federated learning (FL) was introduced by
\citet{mcmahan2017communication} as a privacy-preserving paradigm in
which a large number of clients collaboratively train a shared model
by performing local SGD updates and periodically uploading model
differences to a central server.  This \emph{FedAvg} protocol has
since become the de-facto baseline, yet its convergence degrades
significantly when client data are non-IID—a condition that is the
rule rather than the exception in practice.

\paragraph{Heterogeneity and client drift.}
The root cause of FedAvg's degradation under data heterogeneity is
\emph{client drift}: local objectives diverge from the global
objective, and multiple rounds of local SGD amplify this divergence
before each aggregation step.  \citet{li2018fedprox} address this by
adding a proximal regulariser that keeps local iterates close to the
global model (\emph{FedProx}).  \citet{karimireddy2020scaffold}
provide a more principled fix via control variates (\emph{SCAFFOLD}),
and prove that the resulting algorithm matches the convergence rate of
centralised SGD even under extreme heterogeneity.
\citet{wang2020tackling} further characterise the \emph{objective
inconsistency} that arises when clients take unequal numbers of local
steps, and propose corrective aggregation weights.

\paragraph{Convergence theory.}
Several works have established \emph{linear speedup} in the number of
participating clients.  \citet{yang2021achieving} prove that linear
speedup is achievable for FedAvg under simultaneous non-IID data and
partial client participation, while \citet{qu2020unified} provide a
unified analysis that extends to Nesterov-accelerated variants.
More recently, \citet{jian2025widening} show that the detrimental
effect of data heterogeneity on FedAvg diminishes as the model width
grows, offering a theoretical justification for the empirical success
of over-parameterised networks in FL.

\paragraph{Personalisation.}
An orthogonal strategy is to abandon the goal of a single global
model.  \citet{fallah2020personalized} cast personalisation as a
model-agnostic meta-learning problem (\emph{Per-FedAvg}), while
\citet{dinh2020personalized} use Moreau envelopes to decouple
personalisation from global aggregation.  Although effective, these
methods incur additional communication rounds and store per-client
model copies, increasing both communication and memory overhead—the
very resources we seek to conserve.

\paragraph{Key limitation.}
Despite the richness of the FL literature, the above methods
uniformly assume that clients can afford to compute and store full
gradients and optimiser states.  On memory-constrained edge devices
this assumption breaks down, and no existing FL algorithm
simultaneously addresses convergence under non-IID data
\emph{and} client-side memory constraints.

% ------------------------------------------------------------
\subsection{Communication-Efficient Methods}
% ------------------------------------------------------------

Communication cost is the primary wall-clock bottleneck in FL: each
round requires uploading a full model-sized vector per client.
Two complementary strategies have been developed to reduce this cost:
\emph{gradient compression} and \emph{local training}.

\paragraph{Gradient compression.}
\citet{alistarh2017qsgd} propose QSGD, an unbiased stochastic
quantisation scheme that reduces per-round bandwidth to
$O(\sqrt{d})$ bits while preserving convergence guarantees.
Biased compressors such as top-$k$ sparsification achieve higher
compression ratios but introduce systematic error.
\citet{karimireddy2019error} show that biased compressors must be
paired with an \emph{error-feedback} mechanism to guarantee
convergence, and \citet{richtarik2021ef21} later provide a cleaner
and tighter analysis via the EF21 framework.  The unified EF-BV
theory of \citet{condat2022efbv} reconciles error feedback with
variance reduction for both biased and unbiased compressors.

Fundamental limits have also been established.
\citet{huang2022lower} derive matching lower bounds for distributed
learning with compression, and \citet{he2023unbiased} characterise
precisely when unbiased compression saves communication.
Near-optimal algorithms such as ADIANA
\citep{li2020acceleration} approach these lower bounds by combining
acceleration with compression.

\paragraph{Local training.}
Allowing clients to run multiple local SGD steps before communicating
reduces the number of communication rounds at the cost of increased
client drift.  \citet{basu2020qsparse} combine local SGD with
sparsification and quantisation (Qsparse-local-SGD), while
\citet{tang2019doublesqueeze} apply error compensation to
\emph{both} upload and download directions (DoubleSqueeze).
\citet{condat2024locodl} provide a principled primal-dual framework
(LoCoDL) that provably unifies local training with unbiased
compression, achieving the best known communication complexity for
this setting.

\paragraph{Decentralised settings.}
In the absence of a parameter server, gossip-based protocols
distribute the communication load.
\citet{lian2017can} prove that decentralised parallel SGD (D-PSGD)
can match the convergence rate of centralised SGD, and
\citet{lian2017async} extend this to asynchronous execution.
\citet{koloskova2019decentralized} show that arbitrary communication
compression can be incorporated into decentralised training while
preserving convergence.

\paragraph{Key limitation.}
All of the above methods take the memory footprint of a training step
as given: each client must materialise a full gradient vector and
maintain a full optimiser state (e.g., Adam's first and second
moments).  On devices with limited RAM this is infeasible for large
models.  Reducing communication cost while also reducing memory cost
in a theoretically grounded framework remains an open problem.

% ------------------------------------------------------------
\subsection{Memory-Efficient Training}
\label{subsec:memory}
% ------------------------------------------------------------

Reducing the memory footprint of on-device training has attracted
growing attention, driven by the deployment of large neural networks
on edge hardware.

\paragraph{Low-rank and subspace methods.}
\citet{hu2022lora} introduce LoRA, which freezes pre-trained weights
and injects trainable low-rank decompositions, drastically reducing
the number of active parameters and the associated optimiser states.
\citet{zhao2024galore} extend this idea to full pre-training via
\emph{GaLore}: instead of restricting weight updates, they project
gradients onto a periodically updated low-rank subspace before
passing them to the optimiser.  This reduces the dimension of the
optimiser state from $d$ to $r \ll d$ without modifying the model
architecture.  However, recent analysis reveals that naive subspace
projection under stochastic noise can fail to converge; provably
convergent variants require careful correction of the projection
error \citep{li2023analysis}.

\paragraph{Memory-efficient FL.}
\citet{liu2025fedmuon} apply local matrix orthogonalisation in FL
(\emph{FedMuon}), achieving faster convergence per round.  However,
because client-side orthogonalisation is a non-linear operation, the
standard drift-correction analysis of SCAFFOLD breaks down, and
explicit additional corrections are required.  This observation
underscores a fundamental tension: \emph{memory-efficient local steps
in FL are not drop-in replacements for standard SGD steps}—they
invalidate the aggregation theory that FL convergence proofs rely on.

\paragraph{Activation and state reduction.}
Beyond weight-update compression, gradient checkpointing trades
recomputation for activation memory, and quantisation-aware
approaches such as mixed-precision training reduce state sizes.
\citet{koloskova2019decentralized} and related decentralised methods
also reduce per-node memory indirectly by partitioning the
communication graph.

\paragraph{The gap—motivating our work.}
The survey above reveals three concrete open problems.
%
\begin{enumerate}
  \item \textbf{No unified framework.}  Existing methods address
    memory efficiency (subspace/low-rank training) or communication
    efficiency (compression/local SGD) in \emph{isolation}.  There is
    no algorithm that provably achieves \emph{both simultaneously} in
    heterogeneous FL.  Combining GaLore-style gradient projection with
    SCAFFOLD-style drift correction and EF21-style error feedback
    requires resolving non-trivial interactions among all three
    mechanisms.

  \item \textbf{Missing drift-correction theory.}  When clients
    perform non-Euclidean or subspace-constrained local steps, the
    standard drift-correction analysis breaks down
    \citep{liu2025fedmuon}.  The correct correction mechanism for
    memory-efficient local updates in FL is uncharacterised.

  \item \textbf{Convergence under joint constraints.}  The
    theoretical limits of jointly memory- and
    communication-constrained FL—particularly under non-IID data,
    partial participation, and nonconvex objectives—remain an open
    problem \citep{huang2022lower,he2023unbiased}.
\end{enumerate}
%
Our algorithm and analysis directly address all three gaps, providing
the first provably convergent FL method that achieves joint memory and
communication efficiency under heterogeneous data.
